\documentclass[11pt]{article}
\usepackage{mathblatt}
% gesetzt von werkzeuge/setzer.py (Sorte selbst) aus dem Lernweg MZE
% und den Bankzeilen; Muster bau/proben/2026-10-09/hypotenuse-muster.
\usepackage{enumitem}
\usepackage{adjustbox,varwidth}
\newgeometry{a4paper,margin=18mm,top=15mm,bottom=20mm,footskip=9mm}
\pagestyle{fancy}\fancyhf{}\renewcommand{\headrulewidth}{0pt}
\fancyfoot[L]{\footnotesize\color{mbgrau}MZE \textperiodcentered{} Termwerte berechnen}
\fancyfoot[R]{\footnotesize\color{mbgrau}\thepage}
\setlength{\parskip}{3pt}
\renewcommand{\kreuz}[1]{\mbox{$\square$\ #1}\quad}
\newcommand{\kreuzl}[1]{\par\hangindent1.4em\hangafter1\noindent$\square$\ #1}
\newcommand{\aufg}[3]{\par\noindent\makebox[0pt][r]{#3}\begin{minipage}[t]{\linewidth}#1\end{minipage}\par}
\newcommand{\aufgg}[4]{\par\noindent\makebox[0pt][r]{#4}\begin{minipage}[t]{0.6\linewidth}\vspace{0pt}#1\end{minipage}\hfill
  \begin{minipage}[t]{0.37\linewidth}\vspace{0pt}\centering #2\end{minipage}\par}
\newcommand{\nr}[1]{\makebox[7mm][l]{\textbf{#1}}}
\newcommand{\lsg}[3]{\par\noindent\begin{minipage}[t]{9mm}\textbf{#1}\end{minipage}%
  \begin{minipage}[t]{52mm}\raggedright\bfseries\boldmath #2\end{minipage}\hspace{3mm}%
  \begin{minipage}[t]{\dimexpr\linewidth-64mm\relax}\raggedright\small #3\end{minipage}\par\vspace{4pt}}
\newlength{\kr}
\makeatletter
\newcommand{\karomerk}[2]{\protected@write\@auxout{}{\string\karoseite{#1}{#2}{\thepage}}}
\newcommand{\karoseite}[3]{\expandafter\xdef\csname karo@#1@#2\endcsname{#3}}
\newcommand{\karohinweis}[3]{\karomerk{h}{#1}\ifcsname karo@k@#1\endcsname
  \ifnum\csname karo@k@#1\endcsname=0\csname karo@h@#1\endcsname\relax#2\else#3\fi
  \else#3\fi}
\makeatother
\begin{document}

\noindent{\LARGE\bfseries Termwerte berechnen}\par\smallskip
\noindent{\small\color{mbgrau}Selbstlernheft Klasse 7 \textperiodcentered{} mit Lösungsheft \textperiodcentered{} Kennung MZE}\par\medskip
\noindent\textbf{So nutzt du das Heft:} Lies im Abschnitt das Beispiel. Rechne dann die Aufgaben der Reihe nach und vergleiche jede sofort mit dem Lösungsheft. Kreuze „kann ich“ an, wenn du die letzte Aufgabe (\emph{Ziel}) allein schaffst.\par
\noindent Mehr Aufgaben zu einem Abschnitt: Kennung deinem Nachhilfelehrer schicken (z.\,B. „MZE mehr B“).\par\medskip
{\arrayrulecolor{mbgrau}\renewcommand{\arraystretch}{1.3}
\noindent\begin{tabular}{|>{\bfseries}p{10mm}|p{112mm}|>{\raggedleft\arraybackslash}p{10mm}|>{\centering\arraybackslash}p{14mm}|}\hline
\rowcolor{black!8} & \textbf{Abschnitt} & \textbf{Seite} & \textbf{kann ich}\\\hline
A & Zahl einsetzen und ausrechnen & \pageref{s:A} & $\square$\\\hline
B & Negative Zahlen einsetzen & \pageref{s:B} & $\square$\\\hline
C & Quadrate & \pageref{s:C} & $\square$\\\hline
D & Zwei Variablen und Bruchstrich & \pageref{s:D} & $\square$\\\hline
T & Probetest & \pageref{s:T} & $\square$\\\hline
\end{tabular}}\par\bigskip

\begin{tcolorbox}[colback=mbkasten,colframe=mbgrau,boxrule=0.5pt,arc=1mm,left=3mm,right=3mm,top=2mm,bottom=2mm,title={\bfseries Formel auf einen Blick},coltitle=black,colbacktitle=black!10]
{\Large Klammer $\to$ Hochzahl $\to$ Punkt ($\cdot$\,,\,$:$) $\to$ Strich ($+$\,,\,$-$) von links nach rechts}\par\medskip
In Worten: Malnehmen und Teilen: gleiche Vorzeichen ergeben Plus, verschiedene Minus: $(-2) \cdot (-3) = 6$, \ $(-6) : 2 = -3$.\par\medskip
\end{tcolorbox}
\medskip\noindent\textbf{Typische Fehler – so nicht}\par\smallskip
\noindent\nr{$\times$}$3 \cdot 4 = 12 + 6 = 18$ ist falsch geschrieben, denn $3 \cdot 4$ ist nicht $18$. Richtig: $3 \cdot 4 + 6 = 12 + 6 = 18$.\par
\clearpage\label{s:A}\noindent\colorbox{black!12}{\makebox[10mm]{\Large\bfseries\strut A}}\hspace{3mm}{\Large\bfseries Zahl einsetzen und ausrechnen}\par\smallskip
\noindent Ein \textbf{Term} wie $4x - 3$ enthält eine \textbf{Variable} $x$: einen Platzhalter für eine Zahl. Setzt du für $x$ eine Zahl ein und rechnest aus, erhältst du den \textbf{Termwert}.\par\medskip
\noindent\textbf{Formel}\par\nopagebreak\noindent\hspace*{4mm}\begin{minipage}{\dimexpr\linewidth-4mm\relax}$4x - 3$ \ für \ $x = 5$: \quad $4 \cdot 5 - 3 = 20 - 3 = 17$\end{minipage}\par\medskip
\noindent\textbf{Beispiel}\quad Berechne den Wert des Terms $3x + 2 \cdot (x - 1)$ für $x = 4$.\par\smallskip
\noindent\par\noindent\hangindent6mm\hangafter1\makebox[6mm][l]{\textbf{1}}Für jedes $x$ die $4$ einsetzen; aus $3x$ wird $3 \cdot 4$: \quad $3 \cdot 4 + 2 \cdot (4 - 1)$\par\smallskip\par\noindent\hangindent6mm\hangafter1\makebox[6mm][l]{\textbf{2}}Klammer zuerst: \quad $4 - 1 = 3$, \ also \ $3 \cdot 4 + 2 \cdot 3$\par\smallskip\par\noindent\hangindent6mm\hangafter1\makebox[6mm][l]{\textbf{3}}Punkt vor Strich: \quad $12 + 6 = 18$\par\smallskip\par\noindent\hangindent6mm\hangafter1\makebox[6mm][l]{\textbf{4}}Der Termwert ist $18$.\par\smallskip\par\medskip
\noindent\textbf{Aufgaben}\hspace{1em}{\small\color{mbgrau}\karohinweis{A}{Rechne unten im Karo.}{Rechne auf einem extra Blatt.} Vergleiche jede Aufgabe gleich mit dem Lösungsheft.}\par\smallskip
\aufg{\hangindent7mm\hangafter1\nr{1.}Berechne den Wert von $5x - 4$ für \ a)~$x = 2$ \qquad b)~$x = 3$ \qquad c)~$x = 10$.}{}{}
\medskip
\aufg{\hangindent7mm\hangafter1\nr{2.}Berechne den Wert für $x = 3$. \par\vspace{3pt} a)~$7 + 2x$ \qquad b)~$4 \cdot (x + 2)$ \qquad c)~$20 - 3x$ \qquad d)~$2x + x - 1$}{}{}
\medskip
\aufgg{\hangindent7mm\hangafter1\nr{3.}Berechne die Termwerte von $3x - 2$ und trage sie in die Tabelle ein.}{\adjustbox{max width=\linewidth}{\begin{varwidth}{50cm}\wertetabelle{x}{3x - 2}{1,2,3,4,5}\end{varwidth}}}{}{}
\medskip
\aufg{\hangindent7mm\hangafter1\nr{4.}Ein Fahrradverleih berechnet den Preis in Euro mit dem Term $4 + 3h$ ($h$: Anzahl der Stunden). Tom hat 20\,€ und möchte das Rad 6 Stunden leihen. Reicht sein Geld? Wie viel bleibt übrig oder fehlt?}{}{{\scriptsize\color{mbgrau}Ziel\hspace{2mm}}}
\medskip
\par\vspace{3mm}\setlength{\kr}{\dimexpr\pagegoal-\pagetotal-10mm\relax}%
\ifdim\kr>12mm\karomerk{k}{A}\noindent{\scriptsize\color{mbgrau}Platz zum Rechnen}\par\nointerlineskip\vspace{1mm}\noindent
\begin{tikzpicture}\clip (0,0) rectangle ({\linewidth-0.5mm},\kr);\draw[black!16,line width=0.3pt,step=5mm] (0,0) grid (\linewidth,\kr);\end{tikzpicture}\fi\par
\clearpage\label{s:B}\noindent\colorbox{black!12}{\makebox[10mm]{\Large\bfseries\strut B}}\hspace{3mm}{\Large\bfseries Negative Zahlen einsetzen}\par\smallskip
\noindent Eine negative Zahl setzt du in Klammern ein: Aus $3x$ wird $3 \cdot (-2)$. Dann rechnest du in der Reihenfolge von Seite 1.\par\medskip
\noindent\textbf{Formel}\par\nopagebreak\noindent\hspace*{4mm}\begin{minipage}{\dimexpr\linewidth-4mm\relax}$2x + 9$ \ für \ $x = -3$: \quad $2 \cdot (-3) + 9 = -6 + 9 = 3$ \par\vspace{3pt} $10 - x$ \ für \ $x = -4$: \quad $10 - (-4) = 10 + 4 = 14$ \par\vspace{3pt} $-2x$ \ für \ $x = -3$: \quad $-2 \cdot (-3) = 6$ \par\vspace{3pt} $-3 - 6 = -9$ \ (weiter ins Minus)\end{minipage}\par\medskip
\noindent\textbf{Beispiel}\quad Berechne den Wert des Terms $7 - 3 \cdot (x + 2)$ für $x = -5$.\par\smallskip
\noindent\par\noindent\hangindent6mm\hangafter1\makebox[6mm][l]{\textbf{1}}$-5$ einsetzen: \quad $7 - 3 \cdot (-5 + 2)$\par\smallskip\par\noindent\hangindent6mm\hangafter1\makebox[6mm][l]{\textbf{2}}Klammer zuerst: \quad $-5 + 2 = -3$, \ also \ $7 - 3 \cdot (-3)$\par\smallskip\par\noindent\hangindent6mm\hangafter1\makebox[6mm][l]{\textbf{3}}Punkt vor Strich: \quad $3 \cdot (-3) = -9$, \ also \ $7 - (-9)$\par\smallskip\par\noindent\hangindent6mm\hangafter1\makebox[6mm][l]{\textbf{4}}Minus vor Minus wird Plus: \quad $7 + 9 = 16$\par\smallskip\par\medskip
\noindent\textbf{Aufgaben}\hspace{1em}{\small\color{mbgrau}\karohinweis{B}{Rechne unten im Karo.}{Rechne auf einem extra Blatt.} Vergleiche jede Aufgabe gleich mit dem Lösungsheft.}\par\smallskip
\aufg{\hangindent7mm\hangafter1\nr{1.}Berechne den Wert von $3x + 10$ für \ a)~$x = -2$ \qquad b)~$x = -5$ \qquad c)~$x = -1$.}{}{}
\medskip
\aufg{\hangindent7mm\hangafter1\nr{2.}Berechne den Wert für $x = -3$. \par\vspace{3pt} a)~$8 - x$ \qquad b)~$-2x$ \qquad c)~$5 - 4x$ \qquad d)~$x - 6$}{}{}
\medskip
\aufg{\hangindent7mm\hangafter1\nr{3.}Berechne. \par\vspace{3pt} a)~$6 \cdot (x - 1) + 20$ für $x = -2$ \qquad b)~$10 - 2 \cdot (x + 3)$ für $x = -7$}{}{}
\medskip
\aufg{\hangindent7mm\hangafter1\nr{4.}Welcher Term hat für $x = -2$ den größeren Wert: $3 - 4x$ oder $2 \cdot (x + 6)$? Um wie viel ist er größer?}{}{{\scriptsize\color{mbgrau}Ziel\hspace{2mm}}}
\medskip
\par\vspace{3mm}\setlength{\kr}{\dimexpr\pagegoal-\pagetotal-10mm\relax}%
\ifdim\kr>12mm\karomerk{k}{B}\noindent{\scriptsize\color{mbgrau}Platz zum Rechnen}\par\nointerlineskip\vspace{1mm}\noindent
\begin{tikzpicture}\clip (0,0) rectangle ({\linewidth-0.5mm},\kr);\draw[black!16,line width=0.3pt,step=5mm] (0,0) grid (\linewidth,\kr);\end{tikzpicture}\fi\par
\clearpage\label{s:C}\noindent\colorbox{black!12}{\makebox[10mm]{\Large\bfseries\strut C}}\hspace{3mm}{\Large\bfseries Quadrate}\par\smallskip
\noindent Setz eine negative Zahl auch beim Quadrat in Klammern ein. Dann gilt das Quadrat für die ganze Zahl mit Minus, und es kommt immer eine positive Zahl heraus.\par\medskip
\noindent\textbf{Formel}\par\nopagebreak\noindent\hspace*{4mm}\begin{minipage}{\dimexpr\linewidth-4mm\relax}$(-3)^2 = (-3) \cdot (-3) = 9$ \qquad $x^2 + 1$ \ für \ $x = -3$: \quad $(-3)^2 + 1 = 9 + 1 = 10$\end{minipage}\par\medskip
\noindent\textbf{Beispiel}\quad Berechne den Wert des Terms $2x^2 - 3x$ für $x = -4$.\par\smallskip
\noindent\par\noindent\hangindent6mm\hangafter1\makebox[6mm][l]{\textbf{1}}$-4$ in Klammern einsetzen: \quad $2 \cdot (-4)^2 - 3 \cdot (-4)$\par\smallskip\par\noindent\hangindent6mm\hangafter1\makebox[6mm][l]{\textbf{2}}Hochzahl zuerst: \quad $(-4)^2 = (-4) \cdot (-4) = 16$, \ also \ $2 \cdot 16 - 3 \cdot (-4)$\par\smallskip\par\noindent\hangindent6mm\hangafter1\makebox[6mm][l]{\textbf{3}}Punkt vor Strich: \quad $2 \cdot 16 = 32$ \ und \ $3 \cdot (-4) = -12$, \ also \ $32 - (-12)$\par\smallskip\par\noindent\hangindent6mm\hangafter1\makebox[6mm][l]{\textbf{4}}Minus vor Minus wird Plus: \quad $32 + 12 = 44$\par\smallskip\par\medskip
\noindent\textbf{Aufgaben}\hspace{1em}{\small\color{mbgrau}\karohinweis{C}{Rechne unten im Karo.}{Rechne auf einem extra Blatt.} Vergleiche jede Aufgabe gleich mit dem Lösungsheft.}\par\smallskip
\aufg{\hangindent7mm\hangafter1\nr{1.}Berechne den Wert von $x^2$ für \ a)~$x = 5$ \qquad b)~$x = -5$ \qquad c)~$x = -1$ \qquad d)~$x = -10$.}{}{}
\medskip
\aufg{\hangindent7mm\hangafter1\nr{2.}Berechne den Wert für $x = -2$. \par\vspace{3pt} a)~$4x^2$ \qquad b)~$x^2 + 4$ \qquad c)~$x^2 - x$}{}{}
\medskip
\aufg{\hangindent7mm\hangafter1\nr{3.}Berechne. \par\vspace{3pt} a)~$10 - x^2$ für $x = -3$ \qquad b)~$2x^2 + 5x$ für $x = -1$}{}{}
\medskip
\aufgg{\hangindent7mm\hangafter1\nr{4.}Berechne die Termwerte von $x^2 - 2x$ und trage sie in die Tabelle ein. Für welche Zahl $x$ ist der Termwert negativ?}{\adjustbox{max width=\linewidth}{\begin{varwidth}{50cm}\wertetabelle{x}{x^2 - 2x}{-2,-1,0,1,2,3}\end{varwidth}}}{}{{\scriptsize\color{mbgrau}Ziel\hspace{2mm}}}
\medskip
\par\vspace{3mm}\setlength{\kr}{\dimexpr\pagegoal-\pagetotal-10mm\relax}%
\ifdim\kr>12mm\karomerk{k}{C}\noindent{\scriptsize\color{mbgrau}Platz zum Rechnen}\par\nointerlineskip\vspace{1mm}\noindent
\begin{tikzpicture}\clip (0,0) rectangle ({\linewidth-0.5mm},\kr);\draw[black!16,line width=0.3pt,step=5mm] (0,0) grid (\linewidth,\kr);\end{tikzpicture}\fi\par
\clearpage\label{s:D}\noindent\colorbox{black!12}{\makebox[10mm]{\Large\bfseries\strut D}}\hspace{3mm}{\Large\bfseries Zwei Variablen und Bruchstrich}\par\smallskip
\noindent Jede Variable bekommt ihre eigene Zahl. Ein Bruchstrich wirkt wie eine Klammer um Zähler und Nenner: erst oben und unten ausrechnen, dann teilen.\par\medskip
\noindent\textbf{Formel}\par\nopagebreak\noindent\hspace*{4mm}\begin{minipage}{\dimexpr\linewidth-4mm\relax}$\dfrac{a + b}{c} = (a + b) : c$\end{minipage}\par\medskip
\noindent\textbf{Beispiel}\quad Berechne den Wert des Terms $\dfrac{a + b}{c}$ für $a = 3$, $b = -9$ und $c = -2$.\par\smallskip
\noindent\par\noindent\hangindent6mm\hangafter1\makebox[6mm][l]{\textbf{1}}Jede Variable durch ihre Zahl ersetzen: \quad $\dfrac{3 + (-9)}{-2}$\par\smallskip\par\noindent\hangindent6mm\hangafter1\makebox[6mm][l]{\textbf{2}}Oben und unten zuerst: \quad Zähler $3 + (-9) = -6$, \ Nenner $-2$\par\smallskip\par\noindent\hangindent6mm\hangafter1\makebox[6mm][l]{\textbf{3}}Teilen: \quad $(-6) : (-2) = 3$\par\smallskip\par\medskip
\noindent\textbf{Aufgaben}\hspace{1em}{\small\color{mbgrau}\karohinweis{D}{Rechne unten im Karo.}{Rechne auf einem extra Blatt.} Vergleiche jede Aufgabe gleich mit dem Lösungsheft.}\par\smallskip
\aufg{\hangindent7mm\hangafter1\nr{1.}Berechne für $a = 4$ und $b = -3$. \par\vspace{3pt} a)~$a + b$ \qquad b)~$2a - b$ \qquad c)~$a \cdot b$ \qquad d)~$b - a$}{}{}
\medskip
\aufg{\hangindent7mm\hangafter1\nr{2.}Berechne den Wert des Terms $\dfrac{a - b}{c}$ für $a = 1$, $b = -7$ und $c = -4$.}{}{}
\medskip
\aufg{\hangindent7mm\hangafter1\nr{3.}Berechne den Wert des Terms $\dfrac{a - b}{c}$ für $a = -1$, $b = 7$ und $c = -4$.}{}{}
\medskip
\aufg{\hangindent7mm\hangafter1\nr{4.}Für $a = -2$ und $b = 4$: Welcher Term hat den größten Wert? Kreuze an.\par\medskip\noindent\hspace*{7mm}\mbox{$\square$\ $a \cdot b$}\hspace{10mm}\mbox{$\square$\ $\dfrac{b}{a}$}\hspace{10mm}\mbox{$\square$\ $b - a$}\hspace{10mm}\mbox{$\square$\ $a - b$}\rule[-3mm]{0pt}{8mm}\par\smallskip }{}{{\scriptsize\color{mbgrau}Ziel\hspace{2mm}}}
\medskip
\par\vspace{3mm}\setlength{\kr}{\dimexpr\pagegoal-\pagetotal-10mm\relax}%
\ifdim\kr>12mm\karomerk{k}{D}\noindent{\scriptsize\color{mbgrau}Platz zum Rechnen}\par\nointerlineskip\vspace{1mm}\noindent
\begin{tikzpicture}\clip (0,0) rectangle ({\linewidth-0.5mm},\kr);\draw[black!16,line width=0.3pt,step=5mm] (0,0) grid (\linewidth,\kr);\end{tikzpicture}\fi\par
\clearpage\label{s:T}\noindent\colorbox{black!12}{\makebox[10mm]{\Large\bfseries\strut T}}\hspace{3mm}{\Large\bfseries Probetest}\par\smallskip
\noindent Gemischt wie in der Prüfung, ohne Beispiel und ohne Taschenrechner. Etwa 20 Minuten.\par\medskip
\noindent\textbf{Aufgaben}\hspace{1em}{\small\color{mbgrau}\karohinweis{T}{Rechne unten im Karo.}{Rechne auf einem extra Blatt.} Vergleiche jede Aufgabe gleich mit dem Lösungsheft.}\par\smallskip
\aufg{\hangindent7mm\hangafter1\nr{1.}Berechne den Wert von $6 - 2x$ für $x = -4$.}{}{}
\medskip
\aufg{\hangindent7mm\hangafter1\nr{2.}Welchen Wert hat der Term $4 \cdot (x - 2)$ für $x = -3$? Kreuze an.\par\smallskip\leftskip7mm \kreuz{$-20$}\ \kreuz{$-14$}\par\kreuz{$-4$}\ \kreuz{$20$}}{}{}
\medskip
\aufg{\hangindent7mm\hangafter1\nr{3.}Berechne den Wert von $2x^2 + x$ für $x = -3$.}{}{}
\medskip
\aufg{\hangindent7mm\hangafter1\nr{4.}Berechne den Wert des Terms $\dfrac{a + b}{c}$ für $a = 6$, $b = -2$ und $c = -2$.}{}{}
\medskip
\aufg{\hangindent7mm\hangafter1\nr{5.}Berechne den Wert von $a^2 - b$ für $a = -3$ und $b = -5$.}{}{}
\medskip
\aufg{\hangindent7mm\hangafter1\nr{6.}Fahrschulen rechnen mit dieser Faustformel für den Anhalteweg in Metern ($v$: Geschwindigkeit in km/h): \par\vspace{3pt}\hspace*{15mm}$3 \cdot \dfrac{v}{10} + \left(\dfrac{v}{10}\right)^2$\par\vspace{3pt} Rechne zuerst $\dfrac{v}{10}$ aus. 20\,m vor einem Auto läuft ein Kind auf die Straße. Kommt das Auto vor dem Kind zum Stehen, wenn es a)~30\,km/h fährt, b)~50\,km/h fährt?}{}{{\scriptsize\color{mbgrau}Ziel\hspace{2mm}}}
\medskip
\par\vspace{3mm}\setlength{\kr}{\dimexpr\pagegoal-\pagetotal-10mm\relax}%
\ifdim\kr>12mm\karomerk{k}{T}\noindent{\scriptsize\color{mbgrau}Platz zum Rechnen}\par\nointerlineskip\vspace{1mm}\noindent
\begin{tikzpicture}\clip (0,0) rectangle ({\linewidth-0.5mm},\kr);\draw[black!16,line width=0.3pt,step=5mm] (0,0) grid (\linewidth,\kr);\end{tikzpicture}\fi\par
\end{document}
